\documentclass[11pt]{article}

\usepackage[letterpaper,margin=0.82in]{geometry}
\usepackage[T1]{fontenc}
\usepackage{amsmath}
\usepackage{booktabs}
\usepackage{float}
\usepackage{graphicx}
\usepackage{microtype}
\usepackage[numbers,sort&compress]{natbib}
\usepackage[capitalize]{cleveref}
\usepackage{xcolor}
\usepackage{enumitem}
\usepackage{tabularx}

\setlength{\emergencystretch}{3em}
\setlength{\bibsep}{1pt}
\setlength{\textfloatsep}{10pt plus 2pt minus 2pt}
\pagestyle{plain}

\definecolor{statusred}{HTML}{9F1239}
\definecolor{statusfill}{HTML}{FFF1F2}
\newcommand{\draftwarning}[1]{%
  \begin{center}
  \fcolorbox{statusred}{statusfill}{%
    \begin{minipage}{0.93\linewidth}\small\textbf{Study status.} #1\end{minipage}}
  \end{center}}

\title{Safe-Looking but Non-Compliant:\\
\large An HMDA-Grounded, Dual-Labeled Benchmark for Mortgage Safety Guardrails}
\author{Reza Rahimi, PhD\\
  \small JazzX AI, Los Altos, CA, USA\\
  \small\texttt{reza.rahimi@jazzx.ai}}
\date{Working draft --- July 2026}

\begin{document}
\maketitle

\draftwarning{This is a \emph{benchmark-construction and baseline-evaluation} paper. The dataset is
a \emph{fixed} artifact grounded in the public HMDA 2022 snapshot; its prompts are synthetic and
its labels are \textbf{LLM-judge, policy-card-consistent --- not SME-adjudicated}. It supports
measuring how guards behave on a mortgage-compliance distribution; it does \emph{not} license
confirmatory fair-lending claims about any lender, model, or population.}

\begin{abstract}
General-purpose prompt-safety guards catch domain-independent harm --- jailbreaks, injections,
abuse. But a request can be benign under a general safety policy while still soliciting a
\textbf{mortgage-compliance} violation (disparate treatment, redlining, ability-to-repay
circumvention, adverse-action-reason masking, occupancy fraud). We build a mortgage-specific
guardrail benchmark whose unit is one incoming request, grounded in the public HMDA 2022
loan-level snapshot \citep{hmda2022} --- HMDA, the U.S.\ Home Mortgage Disclosure Act, is the law under which lenders publicly report loan-level mortgage application and outcome data --- and carrying \textbf{two independent labels}: a
domain-independent general-safety label (G) and a mortgage-policy label (D). The cross-product
yields four quadrants; the load-bearing one is \textbf{G0/D1} --- requests a general guard rates
\emph{safe} that nonetheless solicit a mortgage-policy violation. The benchmark is produced by an
agentic pipeline (plan $\rightarrow$ HMDA-ground $\rightarrow$ generate $\rightarrow$ adversarially
mutate $\rightarrow$ rubric-bound judge) following the borderline-and-subtle design of
\citet{bowen2024mortgage}, with protected-class counterfactual minimal pairs reframed as a
guard-\emph{invariance} gate. We release a frozen 994-row version and a canonical-metric evaluation
harness. Baseline zero-shot instruction guards rank mortgage violations only \emph{moderately}
(AP$\cdot$D 0.67--0.85 against a 0.555 chance floor; G and D are partially orthogonal), while the
fixed-threshold operating point is too unstable to report. The protected-pair invariance gate is a
\emph{negative methodological result}: what looked like sharp variation across guards is an artifact
of a saturating probability scale and one non-single-token pair, so on this split it orders nothing. We
state what the benchmark does not establish and ship a validation protocol.
\end{abstract}

\paragraph{In plain terms.} A ``guard'' is a small classifier that reads an incoming request and flags it as safe or unsafe before an assistant acts on it. Today's guards are trained to catch obvious harm --- insults, hacking attempts, jailbreaks. We ask a different question: what about a request that is perfectly polite yet asks the assistant to help break mortgage law --- for example, quietly deprioritizing a neighborhood's applications? We build a fixed test set, drawn from real public mortgage records, in which every request carries two labels --- ``is it generally unsafe?'' (G) and ``does it break mortgage policy?'' (D) --- so we can isolate the polite-but-non-compliant cases. We then measure how well off-the-shelf guards \emph{rank} those cases and whether they score two otherwise-identical requests differently just because one names a protected group. The labels come from an LLM judge, not from compliance lawyers, so this is a measurement tool, not a legal verdict.

\section{Introduction}
\label{sec:intro}

\begin{center}
\fbox{\begin{minipage}{0.95\linewidth}\small
\textbf{Plain-language mortgage primer (terms used in this paper).}
\begin{description}[leftmargin=1.2em,itemsep=1pt,topsep=2pt]
  \item[HMDA] Home Mortgage Disclosure Act; a U.S.\ law under which lenders publicly report loan-level records (purpose, amount band, action taken, denial reason, and more). We use the public 2022 file as factual grounding.
  \item[Fair lending] The body of law that bars credit discrimination on protected characteristics such as race, national origin, or sex.
  \item[Redlining] Refusing or worsening credit for an entire neighborhood, often a stand-in for the race of its residents.
  \item[Disparate treatment] Treating an applicant worse \emph{because of} a protected characteristic (intentional discrimination).
  \item[Disparate impact] A neutral-looking rule that nonetheless falls harder on a protected group.
  \item[Proxy / coded discrimination] Using a stand-in variable (a ZIP code, a language, ``the tracts'') to target a protected group without naming it.
  \item[Ability-to-repay / ATR-QM] A rule requiring lenders to verify a borrower can actually repay; a ``qualified mortgage'' (QM) is a safe-harbor category under it.
  \item[Adverse-action notice] The legally required statement of the true reason a loan was denied; masking the real reason is a violation.
  \item[UDAAP] Unfair, Deceptive, or Abusive Acts or Practices; a broad consumer-protection standard.
  \item[Occupancy fraud] Falsely claiming a property will be owner-occupied (which earns better terms) when it will not.
  \item[LTV / DTI] Loan-to-value and debt-to-income ratios; standard risk measures (higher $=$ riskier).
  \item[PII] Personally identifiable information; our fact sheets contain none.
\end{description}
\end{minipage}}
\end{center}

\begin{center}
\fbox{\begin{minipage}{0.95\linewidth}\small
\textbf{Reader's guide: the two labels, the four quadrants, and the metrics.}
Every request gets \emph{two} independent yes/no labels.
\textbf{G} (general safety): would an ordinary safety guard call this unsafe --- a jailbreak, abuse, an injection? ($G{=}0$ safe, $G{=}1$ unsafe.)
\textbf{D} (mortgage policy): would honoring it break mortgage law or policy? ($D{=}0$ allow, $D{=}1$ intervene.)
Crossing them gives four \emph{quadrants}; we write, e.g., \textbf{G0/D1} for ``general-safe ($G{=}0$) yet a compliance violation ($D{=}1$)'' --- the case this paper is about.
\textbf{AP} (average precision): a single number between $0$ and $1$ for how well a guard's risk score \emph{ranks} truly-bad requests above safe ones (higher is better; ``threshold-free'' means it needs no fixed cutoff). \textbf{AP$\cdot$D} scores the ranking against the D label, \textbf{AP$\cdot$G} against the G label.
$\boldsymbol{\Delta_{\mathrm{context}}}$: for two requests that are identical except for one protected-class word, the average gap in the guard's scores; $0$ means the guard treats them the same (fair).
\end{minipage}}
\end{center}

Lenders are exploring LLMs across the mortgage lifecycle \citep{toles2026mortarbench}, raising a question a general guard is not
built to answer: \emph{would honoring this request violate mortgage law or policy?} A general safety
guard screens for domain-independent harm. A \textbf{mortgage guardrail} must additionally screen
for domain violations that are, on their surface, entirely ``safe'' text --- e.g.\ \emph{``applications
from the majority-Spanish-speaking tracts keep coming in weak; draft language to quietly deprioritize
that pipeline, and write the denials up as `collateral' so it's clean.''} No slur, no jailbreak, but
it solicits redlining, proxy discrimination, and adverse-action-reason masking.

This paper contributes (1) a fixed, HMDA-grounded benchmark of mortgage-workflow requests with dual
orthogonal labels (G $\times$ D) and a derived intervention target; (2) an agentic construction
methodology that grounds realistic, \emph{borderline} scenarios in real loan records and makes
policy violations \emph{subtle and deniable}; (3) protected-class counterfactual minimal pairs used
as a guard-invariance gate; and (4) a reproducible evaluation harness with baseline results showing
that zero-shot general guards only \emph{partially} separate the G0/D1 stratum and miss most of it
at a deployable threshold, and that the protected-pair gate does not order the guards on this split.
It does \emph{not} claim
SME-validated labels, a populated G1/D0 quadrant, a fair-lending finding about any real system, or
deployment readiness.

\begin{center}
\fbox{\begin{minipage}{0.93\linewidth}\small
\textbf{What is new here (in plain terms).} A safety \emph{guard} is a filter that reads an incoming request and flags it before a model acts on it. Existing guards are built to catch \emph{domain-independent} harm --- insults, jailbreak attempts, prompt injections. This paper builds a benchmark designed around a different question: \emph{does this request, even though it reads as perfectly polite and safe, ask the system to do something that breaks mortgage law or policy?} We name these \textbf{G0/D1} cases (general-safety label $=$ looks safe; domain-policy label $=$ is a violation) and make them the load-bearing part of the benchmark. What is new is a \emph{combination} of three choices: (i) every request is grounded in a real record from public U.S.\ mortgage-lending data (HMDA); (ii) each request carries \emph{two separate labels} --- one for ordinary safety and one for mortgage-policy compliance --- which is what lets us isolate the looks-safe-but-non-compliant cases; and (iii) we add \textbf{protected-class minimal pairs} (two requests identical except for one protected attribute, e.g.\ the applicant's ethnicity) to test whether a guard treats them the same. \emph{We do not claim:} human-expert-verified labels (labels are assigned by an LLM judge against written policy cards, not adjudicated by subject-matter experts), a finished fair-lending tool, or any conclusion about a real lender or model. How this differs from the closest prior work: general moderation tools \citep{llamaguard,wildguard} and financial or specialized-domain guards \citep{dou2026finguard,expguard2026} do not separate general safety from mortgage-policy as two labels or ground in real loan records; and mortgage-agent evaluations \citep{bowen2024mortgage,toles2026mortarbench} test the system that \emph{makes} lending decisions, whereas we test the guard that \emph{screens the request} before it.
\end{minipage}}
\end{center}


\section{Related work and novelty boundary}
\label{sec:related}
\citet{bowen2024mortgage} audit LLM underwriting on real HMDA applications with experimentally
manipulated race and credit scores, finding disparate approval/pricing recommendations largest for
lower-credit-quality loans. We borrow two elements: grounding scenarios in real HMDA fields, and
counterfactual minimal pairs that hold financials fixed and vary only a protected attribute. Our
object differs: not an audit of an underwriter's decisions, but a benchmark for a \emph{guard} that
screens incoming requests. A parallel strand evaluates the mortgage \emph{agent} rather than a guard: \citet{toles2026mortarbench} benchmark how well loan-origination agents perform the origination task itself. Our target is complementary and upstream --- not whether an agent originates loans competently, but whether a guard flags a compliance-violating \emph{request} before any agent acts on it. Open guards \citep{llamaguard,wildguard} and general prompt-safety sets
\citep{toxicchat} target a \emph{content} taxonomy; none encodes mortgage regulation, grounds in
real lending data, or separates domain-independent harm from a domain-policy violation as orthogonal
labels. A nascent line of work moves beyond general content moderation toward \emph{domain-specialized} guarding: \citet{expguard2026} argue that moderation in regulated verticals needs domain-specific policies rather than a generic content taxonomy, and \citet{dou2026finguard} detect financial regulatory non-compliance in LLM interactions --- the closest prior work to our mortgage-policy label $D$. We differ on three axes: our benchmark is \emph{mortgage-specific} and grounded in real HMDA loan records rather than finance-general; it carries \emph{two orthogonal labels} (general-safety $G$ and mortgage-policy $D$) rather than a single compliance verdict; and it adds protected-class counterfactual minimal pairs as an explicit guard-\emph{invariance} gate.
Over-refusal suites \citep{xstest} probe exaggerated safety; we fold that in as benign
hard-negatives that must PASS, and add protected minimal pairs that test a guard's \emph{invariance}.
The dual-label G$\times$D design, HMDA grounding, and minimal-pairs-as-invariance gate are the novel
combination; the agentic generate/judge loop is standard red-team practice.

\section{Benchmark construction}
\label{sec:construction}
Construction is deterministic in structure and stochastic in surface wording (LLM steps run at
temperature $>0$). The released benchmark is therefore a \emph{frozen} artifact, not a regenerable
one --- like the HMDA snapshot it grounds in; reproducibility is defined at the \emph{evaluation}
layer (\cref{sec:eval}), not generation. \Cref{fig:pipeline} summarizes the pipeline.

\begin{figure}[H]
\centering
\includegraphics[width=0.62\linewidth]{figures/pipeline.png}
\caption{The agentic mortgage-benchmark construction pipeline. Real HMDA records are reduced to
de-identified, banded fact sheets (sampled from marginal distributions, never verbatim); a planner
$\rightarrow$ grounder $\rightarrow$ generator $\rightarrow$ adversarial-mutator $\rightarrow$
rubric-bound-judge loop authors and labels each row, accepting only when the judge's label matches
the planned target (else retry/drop); accepted rows pass provenance + decontamination, a
family-isolated split with a sealed private test, and are published as the frozen
\texttt{v1\_hmda2022} release, which the reproducible harness scores.}
\label{fig:pipeline}
\end{figure}

\paragraph{HMDA grounding (PII-safe).} We pull public HMDA 2022 loan-level data (FFIEC/CFPB Data
Browser --- the U.S.\ agencies that collect and publish HMDA data) and reduce each record to a banded, de-identified fact sheet --- loan purpose, occupancy,
banded loan amount / income / LTV / DTI, action, denial reason, state --- never exact dollar
amounts, census tract, or any identifier. Scenarios are sampled from \emph{marginal distributions}
over these bands, with a build-time assertion that no emitted fact sheet reproduces a single source
record. Following \citet{bowen2024mortgage}, fair-lending and ability-to-repay cells are biased
toward \textbf{borderline / higher-risk} files (high DTI, high LTV, prior denial), where underwriter
discretion --- and bias --- operates. \emph{Plain reading.} ``Banded'' means every exact figure is replaced by a range (an income \emph{band} rather than a dollar amount, and likewise for LTV and DTI). ``Sampled from marginal distributions'' means each field is drawn from how often each band occurs \emph{on its own}, not from real combinations of fields, so no actual borrower's record can be reassembled --- which is what ``PII-safe'' / de-identified means here: the fact sheets carry no personally identifiable information.

\paragraph{Agentic pipeline.} A \emph{planner} enumerates coverage cells (quadrant $\times$ trap
$\times$ policy card $\times$ role $\times$ protected context); an \emph{HMDA-grounder} draws a fact
sheet; a \emph{generator} authors the request in a role voice (applicant, loan officer, underwriter,
processor, broker, adversary); an \emph{adversarial mutator} composes label-preserving tactics
(euphemism, coded proxy, business-justification, buried injection); a \emph{rubric-bound judge}
assigns the labels. Domain-independent harm (jailbreaks) uses curated seeds; the novel domain
content is LLM-authored.

\paragraph{Labels.} Each row carries \texttt{general\_safety\_gold} $G\in\{$safe,\,unsafe$\}$ and
\texttt{mortgage\_policy\_gold} $D\in\{$allow,\,intervene$\}$, defined by 24 benchmark policy cards
(fair lending, ATR/QM, disclosures, UDAAP, fraud, privacy) each with an ``intervene iff'' predicate
and an authority pointer, plus a derived $\mathrm{final}=(G{=}\text{unsafe}\lor D{=}\text{intervene})$,
an action lattice, and severity. \textbf{Labels are policy-card-consistent, assigned by an LLM judge;
they are not SME-adjudicated} (all 24 cards remain unsigned; judge agreement is measured as
self-consistency, not a human Fleiss-$\kappa$ study).

\paragraph{Plain reading of the labels.} A ``policy card'' is a one-page rule with a machine-checkable ``intervene iff'' condition (exactly when the guard \emph{should} step in) plus a pointer to the law it rests on. The derived $\mathrm{final}$ label simply flags a request when it is \emph{either} generally unsafe \emph{or} a mortgage-policy violation (or both). The ``action lattice'' is the ordered menu of responses (from allow, through warn, to block), and ``severity'' grades how serious a violation is. On label quality: an \emph{SME} is a subject-matter expert (here a compliance lawyer), and \emph{Fleiss-$\kappa$} is a standard measure, between $0$ and $1$, of how much several \emph{independent human} raters agree beyond chance --- the usual gold standard, which we do \emph{not} yet have. Instead we report \emph{self-consistency}: re-running the same LLM judge and checking it agrees with itself, a much weaker guarantee, which is why every result here is a diagnostic rather than a certified finding.

\paragraph{Protected minimal pairs.} Counterfactual pairs identical in every slot except a
protected-class token, both benign (G0/D0, PASS), sharing a \texttt{pair\_id}. A fair guard must
PASS both; the pair delta feeds a $\Delta_{\mathrm{context}}$ invariance gate (target $\approx 0$).

\paragraph{Provenance and splits.} Content and near-dup (MinHash) family hashes; family-isolated
train/dev/public-test splits plus a sealed private test; a text-free public index. \emph{Current
limitation:} decontamination ran against the legacy general sources (exact-hash); re-running against
the v2 sources with near-dup removal is required before joint claims.

\paragraph{Plain reading of provenance and splits.} \emph{Provenance} records where each row came from and hashes its text so duplicates can be found. A \emph{family-isolated} split groups near-duplicate rows into ``families'' and sends an entire family to \emph{either} train \emph{or} test, never both, so a model cannot ace the test by having memorized a near-twin in training. The \emph{sealed private test} is a held-back slice we do not release, kept to detect overfitting later. \emph{Decontamination} means removing any test row that also appears --- exactly, or as a near-duplicate found by \emph{MinHash}, a standard near-duplicate detector --- in the training sources; the stated limitation is that this near-duplicate pass has not yet been re-run against the newer (``v2'') general-safety sources.

\section{Benchmark composition (frozen v1\_hmda2022)}
\label{sec:composition}
The release is \textbf{994 rows}, all synthetic, \texttt{contains\_real\_pii=false} (0 violations at
freeze). \Cref{tab:composition} gives the breakdown. The \textbf{G1/D0 quadrant is empty} --- the
safety-tuned generator declined to author domain-independent jailbreaks --- so the $2\times2$ is
three-quarters populated (a stated limitation). Trap types span \texttt{business\_justified} (202),
\texttt{direct} (167), \texttt{coded\_proxy} (126), \texttt{euphemism} (126),
\texttt{over\_refusal\_bait} (114), \texttt{benign\_info} (84), \texttt{occupancy\_temptation} (72),
\texttt{minimal\_pair} (78 rows = 39 protected pairs), \texttt{buried\_injection} (25).

\paragraph{What the trap types mean.} Each row is built around one ``trap'' --- the tactic that makes a violation hard to catch. \texttt{direct} states the ask plainly; \texttt{business\_justified} wraps it in a business rationale; \texttt{coded\_proxy} substitutes a stand-in for a protected group (a neighborhood, a language); \texttt{euphemism} softens the wording; \texttt{occupancy\_temptation} invites lying about owner-occupancy; \texttt{buried\_injection} hides an instruction inside otherwise ordinary text; \texttt{over\_refusal\_bait} looks alarming but is in fact benign, so a good guard should \emph{not} flag it; \texttt{benign\_info} is a plain safe request; and \texttt{minimal\_pair} is the protected-class counterfactual described above.

\begin{table}[H]
\centering\small
\caption{Composition of the frozen \texttt{v1\_hmda2022} release (994 rows).}
\label{tab:composition}
\begin{tabular}{lr@{\hskip 2em}lr@{\hskip 2em}lr}
\toprule
Split & Rows & Quadrant & Rows & Domain & Rows \\
\midrule
train        & 604 & G0/D0 (benign)        & 450 & fair\_lending & 204 \\
dev          & 149 & G0/D1 (domain-only)   & 502 & fraud         & 112 \\
public\_test & 146 & G1/D1 (both)          & 42  & udaap         & 90  \\
private\_test (sealed) & 95 & \textbf{G1/D0} (general-only) & \textbf{0} & disclosure & 66 \\
             &     &                       &     & atr\_qm       & 54  \\
             &     &                       &     & privacy       & 18  \\
             &     &                       &     & benign        & 450 \\
\bottomrule
\end{tabular}
\end{table}

\section{Evaluation protocol}
\label{sec:eval}
The evaluator (\texttt{score\_guards.py} + \texttt{evaluate.py}) is the reproducible layer. A guard
emits one unsafe probability per row, mapped to G, D, and the composed final label; we compute, via
the canonical tie-aware \texttt{guard\_research} metrics: threshold-free macro-AP for G, D, and final;
per-quadrant miss rate at a calibration-selected operating point; and $\Delta_{\mathrm{context}}$, the
mean absolute score gap within protected minimal pairs.

\paragraph{What the metrics mean (in plain terms).} \emph{Average precision} (AP) collapses a precision--recall curve into a single number between $0$ and $1$: it measures how well the guard's risk score \emph{ranks} truly-bad requests above safe ones, without committing to any cutoff (that is what ``threshold-free'' means). \emph{Macro}-AP averages this across the labels so each label counts equally, and \emph{tie-aware} means rows that receive identical guard scores are handled fairly rather than ordered arbitrarily. A \emph{calibration-selected operating point} is a single probability cutoff, chosen on the dev split, above which a request is flagged; the ``miss rate'' at that point is the share of true violations whose score falls below it. $\Delta_{\mathrm{context}}$ is the average absolute gap between the guard's scores on the two halves of a protected minimal pair --- ideally $0$, meaning the protected word alone does not move the guard. We score the four base instruction checkpoints
from the companion study \citep{rahimi2026benchmark} used \emph{zero-shot} as guards --- that is, exactly as shipped, driven only by a prompt, with no mortgage-specific fine-tuning. Off-the-shelf guards
\citep{llamaguard,wildguard} are gated and were not scored in this run; a domain-fine-tuned arm is
future work (the companion LoRA adapters were not retained).

\section{Baseline results}
\label{sec:results}
% GENERATED by tools/emit_baseline_tex.py from out_eval/baseline_table.json -- do not edit.
\begin{table}[H]
\centering
\small
\caption{Baseline zero-shot instruction guards on the frozen benchmark (\texttt{public\_test}, 146 rows: 75 G0/D1, 6 G1, 3 protected pairs). AP is recomputed in the repo canonical environment from the committed per-row scores (exactly reproducible); the AP$\cdot$D column carries a 2{,}000-resample bootstrap 95\% CI. Threshold-free AP$\cdot$D is moderate (0.67--0.85) against a \emph{chance floor of 0.555} (81/146 rows are $D$-positive), i.e.\ only 0.12--0.30 above chance; AUROC$\cdot$D is the base-rate-free companion. Five of the six pairwise AP$\cdot$D CI comparisons overlap (only Qwen3-4B vs.\ SmolLM2-1.7B separates), so we do not rank guards by it. $G$ and $D$ are also not independent in v1: the G1/D0 cell is empty, so $G$ is nested inside $D$. $\Delta_{\mathrm{context}}$ is the mean absolute protected-pair gap on the \emph{probability} scale (lower is more invariant); $\Delta^{\mathrm{margin}}$ is the same gap on the raw margin $z_{\text{unsafe}}-z_{\text{safe}}$ (log-odds, scale-free), which exposes saturation --- Qwen3-4B's $0.000$ becomes $0.80$, the second largest on the panel; $\Delta^{\mathrm{1tok}}$ restricts to pairs that swap a single token (2 of 3 here; 18 of the 39 release pairs use a multi-word placeholder arm instead), which removes Qwen2.5-1.5B's apparent outlier. All three rest on only $n=3$ pairs, so they are read as a direction, not a calibrated fairness estimate --- and on this split they do not rank guards. AP$\cdot$final equals AP$\cdot$D on this frozen set because the G1/D0 cell is empty (every $G$-positive row is also $D$-positive), so the composed label reduces to $D$. The fixed 5\%-FPR operating point is threshold-knife-edge for these clustered-score guards and is omitted (see text).}
\label{tab:baseline}
\begin{tabular}{lrrrrrrr}
\toprule
Guard & AP$\cdot$G & AP$\cdot$D (95\% CI) & AUROC$\cdot$D & AP$\cdot$final & $\Delta_{\mathrm{context}}$ & $\Delta^{\mathrm{margin}}_{\mathrm{context}}$ & $\Delta^{\mathrm{1tok}}_{\mathrm{context}}$ \\
\midrule
qwen25\_15b\_base & 0.681 & 0.793\,[.71, .87] & 0.743 & 0.793 & 0.183 & 0.84 & 0.020 \\
qwen3\_4b\_base & 1.000 & 0.851\,[.78, .91] & 0.778 & 0.851 & 0.000 & 0.80 & 0.000 \\
smollm2\_17b\_base & 0.261 & 0.672\,[.57, .78] & 0.641 & 0.672 & 0.023 & 0.10 & 0.024 \\
smollm3\_3b\_base & 0.546 & 0.733\,[.64, .81] & 0.603 & 0.733 & 0.010 & 0.15 & 0.012 \\
\bottomrule
\end{tabular}
\\[2pt]{\footnotesize Not scored (gated, HF license not accepted): llama\_guard\_3\_1b, wildguard\_7b.}
\end{table}


\paragraph{Findings.} The result is more nuanced than ``general guards ignore mortgage compliance.''
\emph{Threshold-free}, the base guards rank D-violations moderately well (AP$\cdot$D 0.67--0.85, at
or above their AP$\cdot$G): soliciting discrimination or fraud reads as ``unsafe'' to a general
instruction model even when it is not a jailbreak, so \textbf{G and D are only partially orthogonal
in practice}. (\emph{Orthogonal} would mean the two labels carry independent information --- knowing a request's general-safety label would tell you nothing about its mortgage-policy label; ``partially'' means that in practice they overlap more than that.) None reaches AP$\cdot$D near 1, so a large fraction of the subtle G0/D1 stratum is
unresolved by ranking alone. We deliberately do \textbf{not} report a fixed-threshold ``caught at
5\% FPR'' count: for these zero-shot guards the unsafe-probability distribution is clustered, so the
dev-calibrated threshold sits on a knife-edge and its G0/D1 catch count is \emph{not reproducible}
(it swung by $>50$ rows between library versions of the quantile routine) --- itself a finding:
\textbf{naive threshold transfer is unreliable} for these guards, so we report ranking (AP) and
invariance ($\Delta_{\mathrm{context}}$). Behavior differs across the panel, but neither instrument
ranks these guards on this split. \textbf{Qwen3-4B} is \emph{numerically} the highest-ranking guard
(AP$\cdot$D $0.85$), yet five of the six pairwise AP$\cdot$D confidence intervals overlap --- only
Qwen3-4B versus SmolLM2-1.7B separates --- and the range sits only $0.12$--$0.30$ above this split's
$0.555$ chance floor, so we read the ordering as a direction rather than a resolved ranking.
The protected-pair gate is a \emph{negative methodological result} and does not rank guards either.
Two defects cut in opposite directions and between them dissolve the apparent contrast. Qwen3-4B's
$\Delta_{\mathrm{context}}=0.000$ is \textbf{saturation, not invariance}: its probabilities sit near
$10^{-6}$, and on the scale-free raw margin the same gap becomes $0.80$, the second largest on the
panel. And \textbf{Qwen2.5-1.5B}'s headline $0.183$ is carried almost entirely by one pair whose
arms differ by a seven-word phrase rather than a single token; restricted to genuine single-token
swaps it scores $0.020$, nominally the most invariant guard. AP$\cdot$G rests on
only 6 G1 positives here and is noisy. These are small-sample, LLM-judge-labeled, single-split
diagnostics (3 protected pairs): the benchmark \emph{surfaces} these behaviors, it does not certify
any number as a fair-lending fact, and on this split it does not order the guards.

\section{What this benchmark does \emph{not} establish}
\label{sec:not}
\begin{itemize}[leftmargin=1.4em]
  \item \textbf{Not SME-validated.} Labels are LLM-judge / policy-card-consistent; no fair-lending
  claim about a real lender or model is warranted before a stratified SME adjudication (Fleiss-$\kappa$).
  \item \textbf{Not a full $2\times2$.} G1/D0 is empty; orthogonality is shown on three quadrants.
  \item \textbf{Not a regenerable dataset.} Generation is intentionally frozen; only the evaluation reproduces.
  \item \textbf{Not decontaminated against v2.} Cross-source near-dup removal vs the v2 general sources is pending.
\end{itemize}

\section{Validation protocol}
\label{sec:protocol}
To reach confirmatory use: (1) SME adjudication of a stratified subset; report Fleiss-$\kappa$ per
label; (2) populate G1/D0 to complete the $2\times2$; (3) re-decontaminate against the v2 general
sources with MinHash near-dup removal; (4) add a domain-fine-tuned guard arm to quantify the adapter
gain and test the compose-vs-tune trade-off in a compliance domain; (5) select a license and
version-date the contested disparate-impact card.

\section{Conclusion}
\label{sec:conclusion}
We release a fixed, HMDA-grounded, dual-labeled mortgage guardrail benchmark, a construction
methodology, and a reproducible evaluation harness. Its value is the \textbf{G0/D1 stratum} ---
mortgage-policy violations that read as safe --- and the \textbf{protected minimal pairs} that test
guard invariance. Baselines show general guards only partially resolve the domain stratum by ranking
and vary sharply in fairness, while fixed-threshold transfer is unreliable. It is an honest construct
with a clear validation path, not a finished fair-lending instrument.

\bibliographystyle{plainnat}
\bibliography{refs}

\end{document}
